\documentclass[11pt,a4paper]{article}
\usepackage[margin=1in]{geometry}
\usepackage{amsmath,amssymb,amsthm}
\usepackage{xcolor}
\usepackage{booktabs}
\usepackage{longtable}
\usepackage{enumitem}
\IfFileExists{titlesec.sty}{%
  \usepackage{titlesec}%
  \titleformat{\section}{\large\bfseries\color{harborblue}}{\thesection}{0.7em}{}%
  \titleformat{\subsection}{\normalsize\bfseries}{\thesubsection}{0.6em}{}%
}{}
\usepackage[colorlinks=true,linkcolor=blue!45!black,urlcolor=blue!45!black]{hyperref}
\usepackage{tikz}
\IfFileExists{pgfplots.sty}{%
  \usepackage{pgfplots}%
  \pgfplotsset{compat=1.18}%
}{}
\usetikzlibrary{arrows.meta,positioning,fit,backgrounds,calc,shapes.geometric,shapes.symbols,decorations.pathmorphing,decorations.pathreplacing}
\definecolor{harborblue}{RGB}{30,70,110}
\definecolor{shipred}{RGB}{140,30,30}
\definecolor{seagreen}{RGB}{31,110,70}
\definecolor{slate}{RGB}{71,85,105}
\definecolor{slatedark}{RGB}{30,41,59}
\definecolor{amber}{RGB}{217,119,6}
\definecolor{emerald}{RGB}{16,185,129}
\newtheorem{theorem}{Theorem}
\newtheorem{claim}{Claim}
\theoremstyle{definition}
\newtheorem{change}{Change}
\newtheorem{move}{Move}
\theoremstyle{remark}
\newtheorem*{verdict}{Verdict}
\newtheorem*{adj}{Adjudication}
\setlist{itemsep=1pt,topsep=2pt}

% Shared native-figure styles for the harbor-research corpus (relation-maps,
% regime diagrams, and data plots). Vector TikZ/pgfplots only -- see
% docs/harbor-research/figures/CONVENTION.md for the house rule this follows.
\tikzset{
  relnode/.style={draw,rounded corners,thick,minimum width=2.6cm,minimum height=0.9cm,align=center,font=\small},
  relarrow/.style={-{Stealth[length=2.5mm]},thick},
  regimebox/.style={draw,thick,align=center,font=\small,inner sep=4pt},
}
\IfFileExists{pgfplots.sty}{%
\pgfplotsset{
  harbor curve/.style={
    width=0.46\textwidth,height=5.6cm,
    axis lines=left,
    label style={font=\scriptsize},
    tick label style={font=\scriptsize},
    title style={font=\small,align=center,yshift=2pt},
    legend style={font=\scriptsize,draw=black!25,fill=white,fill opacity=0.92,text opacity=1,rounded corners=1pt,inner sep=3pt},
    thick,
  },
  harbor bar/.style={
    width=0.46\textwidth,height=5.6cm,
    ybar,bar width=6pt,
    ymode=log,log origin=infty,
    axis lines=left,
    label style={font=\scriptsize},
    tick label style={font=\tiny},
    xticklabel style={rotate=40,anchor=east,font=\tiny,inner sep=1pt},
    title style={font=\small,align=center,yshift=2pt},
    legend style={font=\scriptsize,draw=black!25,fill=white,fill opacity=0.92,text opacity=1,rounded corners=1pt,inner sep=3pt},
  },
}
}{}

\title{\textbf{Revising the Treatise}\\[4pt]
\large Granular corrections and directional restructuring for\\ \emph{The Harbor, the Person, and the Economy}\\[3pt]
\normalsize Synthesized from three external reviews and prior analysis --- Document 1 of 4}
\author{Prepared for Erich Owens}
\date{August 2026}
\begin{document}
\maketitle

\noindent\textbf{How to read this.} Three independent reviews now sit on the desk: the \emph{agentsd Blueprint} and the \emph{Port Daddy Analysis v2} (one product-engineering voice, confident and ship-oriented), and \emph{The Harbor After the Harbor} (a rigor voice, counterexample-driven). They agree more than they disagree, and where they disagree the disagreement is itself informative. This document does three things: (\S1) locks the corrections all three reviews converge on, because triple-independent agreement on a defect is the strongest signal a manuscript can get; (\S2) adjudicates the places where the product voice and the rigor voice genuinely conflict, and resolves each; (\S3) proposes the one directional restructuring that the rigor review is right about and that the product review will benefit from commercially. Granular first, grand last.

\section{Corrections with three-way consensus --- make these now}
Each item below was independently flagged by at least two of the three reviews (most by all three) and by the prior analysis. These are not stylistic preferences; they are defects that a hostile referee will find. I state the defect, the fix, and the ricochet --- what else changes once the fix lands.

\begin{change}[The stage game is not a Prisoner's Dilemma; the repeated-game threshold is wrong]
The payoff bimatrix printed in \S VI.7.4 has $A$ preferring $\mathsf{T}$ against $\mathsf{F}$ ($1>0$), so $\mathsf{F}$ is not dominant, $(\mathsf{F},\mathsf{F})$ is not Nash, and the pure equilibria are $(\mathsf{F},\mathsf{T})$ and $(\mathsf{T},\mathsf{F})$. The $\delta^\star\!\approx\!0.253$ trigger-strategy claim and its TLA\textsuperscript{+}/Z3 instance therefore prove nothing about the game as drawn. \textbf{Fix:} either substitute a genuine PD, e.g.\ $(\mathsf{T},\mathsf{T})\!=\!(3,3)$, $(\mathsf{T},\mathsf{F})\!=\!(0,4)$, $(\mathsf{F},\mathsf{T})\!=\!(4,0)$, $(\mathsf{F},\mathsf{F})\!=\!(1,1)$ and re-derive the correct threshold ($\delta \ge 1/3$ for the one-shot-deviation bound with these numbers), or --- better --- stop modelling this as an observed PD at all (see Change 9). \textbf{Ricochet:} the ``claim-signaling proposition'' downstream inherits the corrected threshold; any figure quoting $0.253$ must change.
\end{change}

\begin{change}[Exact-key uniqueness is not mutual exclusion]
Theorem II.6.1 proves at most one live row per \emph{exact normalized key}. It says nothing about overlapping line intervals ($[10,30]$ vs.\ $[20,40]$), a whole-file claim vs.\ a symbol claim, or path ancestry (\texttt{dir/} vs.\ \texttt{dir/file}). \textbf{Fix:} restate the theorem over the conflict predicate you actually enforce --- canonical conflict domains, or a transactional interval/hierarchy overlap test (an interval tree checked inside one \texttt{BEGIN IMMEDIATE}). \textbf{Ricochet:} this is the same object as the SQL-level atomicity the product reviews demand (\texttt{UPDATE \ldots WHERE \ldots RETURNING}); state the exclusion \emph{predicate} once and reuse it for claims and for wallet balances alike. Also add the missing \emph{fencing epoch}: a lease without a downstream broker rejecting stale epochs lets a paused holder wake after expiry and act concurrently with its successor. The registry proves one owner \emph{in the table}, never one \emph{real editor}, until every effect boundary checks the epoch.
\end{change}

\begin{change}[The non-forgeable-identity theorem is false as stated]
Theorem III.6.1 claims a newcomer penalty prevents whitewashing. Counterexample (all three reviews, independently): mature reputation $10$, sanction $2$, newcomer floor $0$ $\Rightarrow$ keeping the sanctioned identity yields $8$, resetting yields $0$; free identities do \emph{not} make the sanction ineffective. The proof's equality $\max(r(i)-\Delta, r_0)=r_0$ holds only when $r(i)-\Delta < r_0$. \textbf{Fix:} state the valid conditional --- whitewashing pays iff a fresh identity's \emph{accessible utility} (credit, access, opportunity) is at least the sanctioned identity's, or iff identities can be \emph{multiplied} to defeat a quorum. Move the critical defense to the \emph{principal}: aggregate reputation, exposure, and sanctions across all identities one principal mints. \textbf{Ricochet:} this is the same principal-aggregation that repairs the Sybil argument (III.6 ex.\ 3), the bond accounting (IV.4), and cross-harbor trust (VII). One correction, four sites healed.
\end{change}

\begin{change}[Signal Detection Theory points the wrong way]
Figure I.6 has evidence-of-danger increasing rightward with ``zoom if $z\ge\beta$,'' so moving $\beta$ left is \emph{more} stringent (more zooms, fewer misses, more false alarms) --- the prose calls that criterion ``lax,'' reversing the direction. And the objective double-counts attention if $C_\text{fa}$ already includes wasted inspection and a separate $a\cdot\mathbb{E}[\text{zooms}]$ term is added. \textbf{Fix:} specify $z\sim p(z\mid \text{danger}, r, \text{verifier})$, inspect iff $z\ge\beta$, $L(\beta)=C_\text{miss}P(\text{miss})+C_\text{inspect}P(\text{inspect})+C_\text{delay}\mathbb{E}[\text{delay}]$, and derive the likelihood-ratio threshold. Say explicitly whether reputation $r$ shifts the prior, the observation model, or the miss cost --- it must not act as a free ``trust scalar,'' and must never defeat hard stakes/irreversibility gates. \textbf{Ricochet:} this is the seed of a real theorem (see Document 3, A2): once $r$ enters as calibrated posterior the regret head is \emph{derived}, not asserted.
\end{change}

\begin{change}[The float plan confuses requester funds with provider collateral]
Protocol IV.4.1 lets the requester fund a ``worker bond'' that is then paid or slashed as if it were the worker's money. \textbf{Fix:} separate four buckets --- requester bounty escrow, provider performance bond, optional insurer reserve, fees --- each signed by both principals, with returned collateral explicitly \emph{not} income. A partial settlement must name which evidence units were accepted and whose posted asset moves. \textbf{Ricochet:} amendments now require both principals to sign a new plan version and atomically re-fund each bucket; this is also the hook the escrow/underwriting chapter needs later.
\end{change}

\begin{change}[Anchor attenuation must compare every edge, not every node to the root]
Algorithm V.4 compares each $\mathrm{cap}_i$ to $\mathrm{cap}_\text{root}$, which admits $\texttt{write}\!\to\!\texttt{read}\!\to\!\texttt{write}$: the final write is within root even though it expands the middle node's authority. The appendix finds the attack; the main algorithm and status table still assert the false property. \textbf{Fix:} verify $\mathrm{cap}_i \subseteq \mathrm{cap}_{i-1}$ at every hop (action, resource, audience, conditions, quotas, and $\text{expiry}_i \le \text{expiry}_{i-1}$), bind subjects to keys with proof of possession, and add unique token IDs, ancestor revocation, epoch/freshness, and replay handling. \textbf{Ricochet:} the Kani harnesses must then prove a \emph{relational} property quantified over symbolic parent/child capabilities --- calling a comparator on two concrete vectors is not a universal proof, and the status table must be down-graded to say so (see Change 8).
\end{change}

\begin{change}[Schema-by-union is a latent corruption bug; WAL auto-checkpoint does not bound loss time]
Two systems defects with the same root (the single writer must own \emph{all} global facts): (a) modules self-initializing via \texttt{CREATE TABLE IF NOT EXISTS} let a later module under old code lazily re-add a renamed column, splitting writes across versions --- replace with a linear migration ledger keyed on \texttt{PRAGMA user\_version} plus a schema-epoch/compatibility table; (b) SQLite's WAL auto-checkpoint is a \emph{page-count} threshold, not a timer, so it cannot bound elapsed vulnerability $\Delta t$ --- a quiet database may leave a commit uncheckpointed indefinitely. \textbf{Fix for (b):} expose durability as a per-transaction type \{\textsc{process\_crash} $\mid$ \textsc{power\_loss}\}; route money/authority/non-repeatable-effect writes through \texttt{synchronous=FULL} and verify commit before returning; keep ordinary claims on \texttt{NORMAL}. Only with an independently guaranteed minimum write rate $w_\text{min}$ can you claim $\Delta t \le P/w_\text{min}$, and the manuscript states no such assumption.
\end{change}

\begin{change}[Cryptography cannot prove reality; ``reference monitor'' overstates the same-user design; Kani proves less than claimed]
Three honesty corrections that share a spine. (a) Signatures, hashes, Merkle inclusion, and append-only logs establish attribution and tamper-evidence \emph{relative to a commitment}; they do not establish that every event was logged, that a description is true, that a test measures the intended property, that one party caused a business outcome, or that no omitted channel exists. Introduce typed evidence channels (deterministic verifier, provenance witness, protected meter, expert judgment, buyer preference, delayed outcome) and make every claim name the channel that warrants it and what that channel cannot show. (b) A reference monitor must be tamper-resistant, always-invoked, and small enough to verify; a same-user daemon with advisory hooks and bypassable git/fs/net paths is a \emph{coordinator}, not a complete monitor. Replace the blanket term with the five assurance levels of Change 12. (c) A Kani harness that stubs signature verification and asserts absence of panic proves bounded memory-safety, not the prose security property; state each harness's preconditions, bounds, and postcondition, prove parser refinement separately from assumed crypto, and use a relational constant-time checker for timing. \textbf{Ricochet:} the maturity ledger --- your best feature --- becomes \emph{more} credible once these three claims are scoped, not less.
\end{change}

\section{Where the reviews conflict --- adjudications}
The interesting work. On these points the product voice and the rigor voice pull in opposite directions. Each adjudication states both positions fairly, then resolves.

\begin{adj}\textbf{The core reframe: ``economy of persons'' vs.\ ``evidence-bearing work unit.''} The rigor review's central claim is that the volume keeps almost discovering a stronger thesis --- an epistemic and effect-control system for accountable work --- then covering it with personhood, reputation, and market machinery; it proposes demoting personhood from foundation to optional policy and making the primitive a durable work unit $W = \langle$principal, actor-lineage, intent, acceptance-policy, authority-envelope, expected-effects, claims, evidence-obligations, checkpoint, risk-guarantee, outcome$\rangle$. The product review does the opposite: it leans \emph{into} personhood, feudalism, ``knights wearing colors,'' as go-to-market narrative. \emph{Resolution:} these conflict only if you conflate the formal spine with the marketing surface. The rigor review is correct about the \textbf{formal ontology} --- the work unit is the better primitive, personhood is an inheritance rule (``which liabilities, permissions, evidence, and commitments does a successor inherit?'') and not a claim about a soul. The product review is correct that \textbf{feudal metaphor sells} to the enterprise buyer. Adopt the work unit as Chapter 1's primitive; keep the personhood language as an explicitly-labelled \emph{interpretive interlude} and as external marketing copy, never as a critical premise. This is the directional move of \S3.
\end{adj}

\begin{adj}\textbf{Checkpoints: ``the JSON ledger \emph{is} the neural state'' vs.\ provider-portable checkpoints are impossible.} The product review claims event-sourcing the exact JSON context array plus a Git SHA, replayed under prefix caching, restores neural state for ``fractions of a cent'' --- ``the JSON ledger IS the neural state.'' The rigor review gives an impossibility argument: two hidden provider states $h_1\neq h_2$ can export the same transcript and capsule $C$ yet require different next actions, so a provider-neutral restorer seeing only $C$ cannot preserve both trajectories. \emph{Resolution:} the rigor review is correct and the product review is overclaiming, but the \emph{mechanism} the product review describes is exactly the right \emph{implementation} of the rigor review's ``portable task capsule.'' Reconcile precisely: prefix-cache replay of the message array restores behavior \textbf{within one provider/model/version} (where the transcript plus sampling seed is very nearly sufficient), and is genuinely cheap; it does \textbf{not} restore behavior across providers, because KV-cache, sampling state, and hidden activations are not exported. Keep the event-sourced ledger as the capsule substrate; replace ``IS the neural state'' with ``is a sufficient statistic for same-provider continuation and a semantic capsule for cross-provider \emph{succession}.'' Define correctness at the work-protocol boundary (can the successor close what the ancestor opened?), not bit-identity. The ``big resume button'' survives intact --- it is \emph{task} continuity with \emph{actor} replacement, which is a better product story anyway because it is provider-independent.
\end{adj}

\begin{adj}\textbf{The B2B airlock: ``silicon-enforced NDA, mathematically cannot phone home'' vs.\ the release gate is the real boundary and $q\cdot b$ bits leak.} The product review sells dynamic taint analysis: tag a confidential-reading agent \texttt{[TAINT:HIGH]}, drop egress packets on any tainted network attempt, and Bob's data ``mathematically cannot'' leave. The rigor review is blunt that token-level taint through an LLM is not soundly definable (the model taints everything it writes with everything it read), that an ordinary remote model API is \emph{itself} a trusted reader of plaintext, and that any Erin-visible channel of $b$ freely-chosen bits per job leaks up to $b$ bits, $q\cdot b$ across $q$ jobs, before timing channels. \emph{Resolution:} the rigor review is decisively correct and the product framing is dangerously overconfident --- shipping ``mathematically cannot phone home'' as a security claim invites a breach and a lawsuit. Adopt the rigor architecture verbatim: control the \emph{channel}, not the token (conservatively taint the whole worker after first secret read); make the only exits two declassification gates; price each gate as an information-theoretic leakage budget (fixed schemas, padding, coarse statuses, scheduled release, differential-privacy accounting on aggregates); run the model \emph{inside} the confidential domain because an external API is another reader. The taint-drop-egress mechanism is still useful as a \emph{containment} layer beneath the gate --- it is the ``confine'' rung of the risk ladder --- but it is not the guarantee. This is exactly the clean-room design in Document 3 (C1: Rushby noninterference modulo declassification) and it makes the far stronger honest claim: \emph{every allowed release is explicit and bounded}.
\end{adj}

\begin{adj}\textbf{Rate-the-raters: model heterogeneity + VRF honeypots vs.\ $\rho^\star=G/(dB)$ plus grounded termination.} The product review ``solves'' Conjecture III.11.1 by mandating a Claude/Llama/GPT judge mix (to break correlated mode collapse) and injecting VRF-selected honeypot tasks so a judge never knows if it is being audited. The rigor review says the honest local condition is $\rho\, d\, B > G$, giving critical rate $\rho^\star = G/(dB)$, and that a \emph{recursive} tower needs either finite appeal depth with contractual finality or an independently grounded oracle plus a per-level contraction $\lambda<1$. \emph{Resolution:} these are not in conflict --- they are two halves of one theorem, and neither review states the whole. Heterogeneity is precisely how you obtain the \emph{disjoint cliques} that a contraction proof needs (see Document 3, B2: sealed sampling from $q$ independent cliques gives $\lambda = O(1/q)$); VRF honeypots are the concrete \emph{randomized audit} realizing rate $\rho^\star$; and $\rho^\star = G/(dB)$ is the per-level rate. Fuse them: the corrected Chapter III should present the inspection-game stage result ($\rho^\star$), the contraction under heterogeneous sealed sampling ($\lambda$), and honeypots-as-audit as the implementation, with the honest terminal condition (an exogenously honest root, the operator at $n{=}1$, or a bonded arbitration market at scale). This converts a conjecture into a parameterized theorem with a named implementation --- the best possible outcome for that chapter.
\end{adj}

\begin{adj}\textbf{Inception Canaries: ship it (product) vs.\ it is operator deception (implied by the rigor frame).} The product review proposes having the sidecar secretly inject a subtle vulnerability into a benign PR and score the operator on whether they catch it, gamifying vigilance with a ``God Mode'' unlock. The rigor review never addresses it, but its entire stance --- claims must carry warrant, evidence channels must be honest, the operator is a principal owed legible reasons --- is in tension with a mechanism whose premise is \emph{deceiving your own operator} and, worse, \emph{deliberately writing a real vulnerability} that could escape the canary harness. \emph{Resolution:} keep the intent (combat the abdication gradient; measure and reward vigilance) and change the mechanism to remove deception and live-vulnerability risk. Use \textbf{consented, labelled} audit probes drawn from a blinded, continuously-refreshed holdout of \emph{historical, already-fixed} defects replayed in a sandbox --- never a novel injected vulnerability in a live PR, and never without the operator knowing that probes exist (they need not know \emph{which} item is a probe). This preserves the calibration signal and the dopamine loop while (a) never creating a real vulnerability that can leak past the harness, and (b) not training operators to distrust the very system meant to earn their trust. The prior analysis's Campbell's-law treatment (canaries decay once their signatures are learned; separate the canary-generator from the operator; correct for exposure time) applies directly.
\end{adj}

\begin{adj}\textbf{Tone: ``cull the honest-label fatigue, state with absolute confidence'' vs.\ the honesty \emph{is} the contribution.} Both product reviews want the relentless ``we do not overclaim'' apologies stripped and functionality stated with absolute confidence, caveats relegated to end-of-chapter ``Limitations \& Boundaries'' sections. The rigor review, by contrast, treats the maturity ledger (implemented/partial/specified/proposed) as arguably the \emph{meta-contribution} of the volume. \emph{Resolution:} both are right about different things. The product reviews are right that \emph{prose-level} apologetic tics (``honestly not built yet,'' repeated mid-sentence) induce fatigue and read as fragility --- cut those. The rigor review is right that the \emph{structured} maturity labels and the explicit boundary sections are the volume's signature and must stay. So: keep the four-valued status tags and add a clinical ``Limitations \& Boundaries'' section per chapter; delete the in-prose apologies. State what the system \emph{does} with confidence; state what it does \emph{not} do with precision; stop \emph{apologizing} for the gap between them. The distinction is between honest \emph{structure} (keep, it is rare and valuable) and anxious \emph{tone} (cut).
\end{adj}

\section{The one directional restructuring}
Beyond the line-item fixes, the rigor review makes a structural argument the manuscript should accept, and which --- read correctly --- \emph{strengthens} the product thesis rather than diluting it.

\begin{move}[Re-spine the volume around the work unit; keep the seven papers as essays]
Reorganize the eight-chapter spine so the causal dependency is honest --- \emph{first define and secure the work, then preserve and share it, only then price it}:
\begin{enumerate}[leftmargin=1.5em]
\item \textbf{The Evidence-Bearing Work Unit} --- the primitive, its lifecycle (proposed $\to$ authorized $\to$ executing $\to$ witnessed $\to$ verified\,$\mid$\,contested\,$\mid$\,abandoned), and the receipt.
\item \textbf{The Epistemic Kernel} --- claims, evidence, why/where/why-not provenance, compaction, attention-as-control.
\item \textbf{The Effect Hypervisor} --- authority, roles, confinement, operation safety, and the five assurance levels.
\item \textbf{Continuity Without Metaphysics} --- task capsules, actor replacement, inheritance policy (this is where Parfit becomes an interlude, not a premise).
\item \textbf{The Project Institution} --- people, agents, user needs, typed conflicts, prospective assurance (Alice/Bob/Carl).
\item \textbf{The Sealed Harbor} --- mutual confidentiality, information-flow control, attestation, declassification (Derek/Erin).
\item \textbf{Underwriting Autonomous Work} --- principals, receipts, bonds, adjudication, correlated risk (the economy, correctly subordinated).
\item \textbf{Federation} --- what may gossip (monotone facts, via CALM/invariant-confluence), what needs an authority point (revocation, uniqueness, settlement), and how local authorization crosses boundaries.
\end{enumerate}
The seven original papers survive as essays or interludes. \textbf{Why this helps the product, not just the paper:} the enterprise wedge (compliance, liability containment, the airlock) sells the \emph{work receipt} and the \emph{effect hypervisor} --- chapters 1 and 3 --- not the sovereign economy. Leading with the work unit puts the commercially critical idea first and defers the economy (chapter 7) to where it belongs: after there are real receipts and real failure rates to price. The product reviews implicitly agree --- their Stage 1 wedge is exactly ``API firewall + liability containment,'' i.e.\ the effect hypervisor and the receipt, with the market last.
\end{move}

\noindent\textbf{What to cut.} The rigor review estimates the volume could lose roughly forty percent without losing a core idea: repeated reader-orientation maps, the estimator catalogs, tutorial-level token material, ornamental theorem labels, and the recurring Hobbes/Locke/Parfit/Sen detours (excellent sources of questions, poor certifiers of systems guarantees). I concur, with one refinement from the prior analysis: keep \emph{one} philosophical frame per chapter as an explicitly-labelled interlude (Elster's precommitment for the operator-override chapter earns its place; Hobbes for the state-of-nature framing sells the wedge), and move the rest to endnotes. The maturity ledger and the boundary sections are the exception --- they expand.

\vspace{1em}
\noindent\hrulefill

\noindent\emph{Directional summary.} Fix the eight consensus defects (they are real and a referee will find them). Adopt the rigor review's work-unit spine and the channel-level clean-room; adopt the product review's event-sourcing mechanism, escrow design, and enterprise wedge --- but strip the two overclaims (``JSON is the neural state,'' ``mathematically cannot phone home'') that a careful reader or an adversary would puncture. Keep the honesty \emph{structure}, cut the honesty \emph{tone}. Lead with the work unit; price it last. Documents 2--4 build the product, the proofs, and the papers on this corrected base.
\end{document}
